%%%%%%%%%%%%%%%%%%%%%%%%%%%%%%%%%%%%%%%%%%%%%%%%%%%%%%%%%%%%%%%%%%%%%%%%%%
%   FEHIME CEREN AY - SØKNADSBREV
%   Stilling: Dataanalytiker, DNB Bank ASA (FINN-kode 468593513)
%   Søknadsfrist: 2. august 2026
%%%%%%%%%%%%%%%%%%%%%%%%%%%%%%%%%%%%%%%%%%%%%%%%%%%%%%%%%%%%%%%%%%%%%%%%%%

\documentclass[10pt,a4paper,roman]{moderncv}

\moderncvstyle{casual}
\moderncvcolor{blue}

\usepackage[utf8]{inputenc}
\usepackage[norsk]{babel}
\usepackage[scale=0.78,top=1.6cm,bottom=1.4cm]{geometry}
\setlength{\parskip}{0.9em}

%----------------------------------------------------------------------------------------
%   AVSENDER
%----------------------------------------------------------------------------------------

\firstname{Fehime Ceren}
\familyname{Ay}
\address{Bestumveien 37C, 0281}{Oslo}
\phone{+47 92 25 66 45}
\email{fcerenay@gmail.com}
\extrainfo{\href{https://cerenay.github.io}{\color{blue}cerenay.github.io}}

%----------------------------------------------------------------------------------------
%   MOTTAKER
%----------------------------------------------------------------------------------------

\recipient{Thomas Michael Cassidy Svennevik, avdelingsleder}{DNB Bank ASA \\ Dronning Eufemias gate 30 \\ 0191 Oslo}
\date{30. juli 2026}
\opening{Søknad på stilling som dataanalytiker}
\closing{Med vennlig hilsen,}
\enclosure[Vedlegg]{CV}

\begin{document}
\makelettertitle

Jeg søker med dette på stillingen som dataanalytiker i DNB. Jeg er dataanalytiker
med doktorgrad i økonomi og fem års erfaring med å bruke data til å svare på
spørsmål virksomheter faktisk trenger svar på, først som research scientist i
Telenor Research og siden 2023 som seniorrådgiver i Skatteetatens analysemiljø.

Det som gjør denne stillingen særlig interessant for meg, er at dere er midt i en
modernisering der fagområdene selv skal eie dataproduktene sine. Jeg har jobbet i
to organisasjoner med samme utfordring: dataene finnes, men ikke alltid i en form
som gjør dem anvendbare, og sentrale måltall betyr ofte litt forskjellige ting
avhengig av hvem du spør. Mye av jobben min har handlet om å rydde i nettopp
dette: å bli kjent med både forretningen og dataene godt nok til å kunne
definere begreper som holder, og å legge analysene i en form andre kan bygge
videre på.

Et konkret eksempel: i år publiserte jeg og tre kolleger en analyse av hvorfor
noen Airbnb-utleiere ikke rapporterer leieinntektene sine. Vi koblet
tredjepartsdata fra utleieplattformene mot innleverte skattemeldinger og en
spørreundersøkelse til drøyt 18 000 personer, og brukte regresjons- og
segmenteringsanalyse til å skille mellom fire grupper med helt ulike årsaker til
manglende rapportering. Poenget med analysen var ikke å beskrive et gap, men å gi
etaten et grunnlag for å bruke ressursene der de faktisk virker. Det er den typen
arbeid jeg vil gjøre mer av: starte i en uklar problemstilling, finne og
kvalitetssikre dataene som skal til, og lande på noe som er konkret nok til å
handle på.

Fra Telenor har jeg erfaring med analyse på store datamengder på tvers av
forretningsenheter i Norden og Asia, blant annet prediktive modeller og
tekstanalyse av kundeservicelogger for å måle kundetilfredshet. Der lærte jeg
også hvor mye som avhenger av at analysen kommuniseres riktig. De samme funnene
måtte presenteres på én måte for produktteamene og en annen for konsernledelsen.

Til daglig jobber jeg i SQL, Python og R på Databricks, og jeg bygger
selvbetjente løsninger i Power BI og R Shiny slik at fagmiljøene kan følge opp
egne tall uten å gå veien om meg hver gang. Jeg samarbeider tett med data
engineers, IT og AI-miljøer, og jeg gir gjerne tilbakemeldinger oppstrøms når
datakvalitet eller struktur står i veien for god analyse. Jeg jobber like godt på
norsk som på engelsk.

Doktorgraden min er i økonomi, med vekt på eksperimentelle og kvantitative
metoder. Den har gitt meg metodisk trygghet og vane med å være ærlig om hva
dataene faktisk viser, og hva de ikke viser. Samtidig har de siste fem årene
handlet om anvendt analysearbeid i drift, med de prioriteringene og fristene som
følger med.

Jeg stiller gjerne til samtale, og svarer på det dere måtte lure på.

\makeletterclosing

\end{document}
